% Generated draft: requires mathematical and source review.
\documentclass[11pt]{article}
\usepackage[margin=25mm]{geometry}
\usepackage{fontspec}
\setmainfont{Latin Modern Roman}
\usepackage{amsmath,amssymb,mathtools,unicode-math}
\setmathfont{Latin Modern Math}
\usepackage{enumitem,needspace,xurl,hyperref}
\setlength{\parindent}{0pt}
\setlength{\parskip}{5pt}
\setlength{\emergencystretch}{3em}
\newcommand{\N}{\mathbb{N}}
\newcommand{\Z}{\mathbb{Z}}
\newcommand{\Q}{\mathbb{Q}}
\newcommand{\R}{\mathbb{R}}
\newcommand{\C}{\mathbb{C}}
\newcommand{\F}{\mathbb{F}}
\newcommand{\A}{\mathbb{A}}
\newcommand{\PP}{\mathbb{P}}
\newcommand{\E}{\mathbb{E}}
\newcommand{\eps}{\varepsilon}
\newcommand{\dd}{\,\mathrm{d}}
\begin{document}
\section*{UnsolvedMath ID: 11000018}
\textbf{Source problem code:} catalogue entry
\subsection{Definitions and mathematical statement}
For the usual setting of closed Riemann surfaces \(X_g\) of genus \(g\geq2\), write \(\operatorname{Aut}(X_g)\) for the finite group of holomorphic automorphisms. Fix a property \(P\) of finite groups, such as nilpotence, solvability or being a \(p\)-group. The requirement concerns the \emph{full} automorphism group, not merely a subgroup acting on the surface.

The problem asks for a version of Larsen's counting theorem for the genera at which the relevant upper automorphism-order bound is attained by some \(X_g\) with \(\operatorname{Aut}(X_g)\) having \(P\); it also asks for analogous lower-bound counts and the least genus attaining each given bound. To expose the counting quantifiers, if \(U_P(g)\) denotes a specified applicable upper bound, the relevant counting function is
\[
 \#\{g:2\leq g\leq T,\ \exists X_g\text{ with }
 P(\operatorname{Aut}(X_g)),\ |\operatorname{Aut}(X_g)|=U_P(g)\}.
\]
The bound \(U_P\), Larsen's theorem and its precise asymptotic form are not supplied. For lower bounds, the source also does not say whether it bounds individual orders or the maximum attainable order
\[
 A_P(g)=\max\{|\operatorname{Aut}(X_g)|:P(\operatorname{Aut}(X_g))\},
\]
when this set is nonempty. Attainment of a lower bound for \(A_P(g)\) must not be replaced by mere existence of one surface of that order. Least-genus requests refer to the nonempty attainment sets once the intended bounds are fixed.

The continuation notes that many extremal examples are arithmetic, meaning quotients of \(\mathbb H^2\) by arithmetic lattices, and describes them as highly symmetric but not typical. This motivates a following question about typical surfaces with symmetry, whose statement is not included; no probability measure or typicality theorem is supplied.
\subsection{Short English statement}
For finite-group properties such as nilpotence, solvability or a \(p\)-group condition, count genera attaining the relevant automorphism-order bounds and find the first attaining genus, with the full automorphism group required to have the property.
\subsection{Sources}
\begin{itemize}
\item \emph{Problems on Mapping Class Groups and Related Topics}, edited by Farb, Problem 2.16, ``Automorphism groups with special properties,'' supplied PDF page 29. \url{https://www.math.uchicago.edu/~farb/papers/mcgbook.pdf}. Recorded access 2026-07-29; extraction pdf-font-structure. Larsen's theorem and preceding bounds are referenced but absent.
\item The supplied research summary mentions Zomorrodian's nilpotent bound without its value, hypotheses or complete bibliography entry.
\item Local UnsolvedMath registry, \texttt{data/unsolvedmath/problems.json}, permanent ID \(11000018\), exact source code \texttt{catalogue entry}; supplied OPEN-TRIAGE summary reports no general frequency/sharpness/least-genus resolution.
\end{itemize}
\end{document}
